% =====================================================================
%  setup/diagrams.tex -- TikZ / pgfplots / forest set-up and the shared
%  diagram styles. Every diagram in the book should be built from these.
%
%  Figure pattern (non-floating, stays where it is written):
%    \begin{figure}[H]
%      \centering
%      \begin{tikzpicture}[diagram] ... \end{tikzpicture}
%      \caption{What the reader should notice.}\label{fig:c01-depth-height}
%    \end{figure}
%  In chapter openers / chapter ends use \caption*{...} (no number).
%
%  General styles (shared with the maths book):
%    diagram            picture-level defaults (put in [ ] of tikzpicture)
%    numnode            circled number            boxnode   rounded box with text
%    cellnode           square table/grid cell    hl, hlblue, hlgreen, hlred, hlpurple  fills
%    flow               main arrow (blue, thick)  softarrow secondary arrow
%    guide              dashed helper line        brace / mbrace  curly brace (mirror)
%    annot              small italic annotation
%
%  Tree styles (manual TikZ, for annotated trees):
%    tnode              tree node: circle, blue border, light fill
%    tnodeA/troot       accented node variants (orange = root, green = leaf)
%    tleaf              leaf node (green border)
%    nullnode           small open circle for a missing child
%    tedge              parent-to-child edge (thick dark line)
%    tedgehl            highlighted edge (orange)
%    subtreebox         dashed fit box around a subtree
%
%  forest styles (quick clean trees): \begin{forest} bintree [A[B][C]] \end{forest}
%    bintree            binary tree, symmetric angles, circle nodes
%    genetree           general (n-ary) tree
%    valbintree         binary tree with square value nodes
% =====================================================================
\usetikzlibrary{arrows.meta, positioning, calc, fit, backgrounds, matrix,
  shapes.geometric, shapes.multipart, decorations.pathreplacing,
  decorations.pathmorphing, patterns, shadows, chains, math}
\pgfplotsset{compat=1.18}

\tikzset{
  >={Stealth[length=2.2mm]},
  diagram/.style={font=\small, line cap=round, line join=round},
  numnode/.style={circle, draw=primaryblue, fill=primaryblue!8, thick,
    minimum size=7mm, inner sep=1pt, font=\small},
  boxnode/.style={rectangle, rounded corners=2pt, draw=darkgray, fill=white, thick,
    minimum height=6mm, inner sep=3pt, font=\small, align=center},
  cellnode/.style={rectangle, draw=darkgray!60, minimum size=6mm, inner sep=0pt,
    font=\footnotesize},
  hl/.style={fill=accentorange!28},
  hlblue/.style={fill=primaryblue!18},
  hlgreen/.style={fill=successgreen!22},
  hlred/.style={fill=warningred!18},
  hlpurple/.style={fill=softpurple!20},
  flow/.style={->, thick, primaryblue},
  softarrow/.style={->, darkgray!75},
  guide/.style={dashed, darkgray!55},
  brace/.style={decorate, decoration={brace, amplitude=5pt, raise=2pt}, thick, darkgray},
  mbrace/.style={decorate, decoration={brace, mirror, amplitude=5pt, raise=2pt}, thick, darkgray},
  annot/.style={font=\footnotesize\itshape, text=darkgray, align=center},
  % ---- tree nodes and edges (manual TikZ) ----
  tnode/.style={circle, draw=primaryblue, fill=primaryblue!8, thick,
    minimum size=8mm, inner sep=0pt, font=\small},
  troot/.style={tnode, draw=accentorange, fill=accentorange!18},
  tleaf/.style={tnode, draw=successgreen, fill=successgreen!16},
  nullnode/.style={circle, draw=darkgray!50, fill=white, minimum size=2.2mm,
    inner sep=0pt},
  tedge/.style={-, thick, darkgray!75},
  tedgehl/.style={-, very thick, accentorange},
  subtreebox/.style={draw=softpurple, dashed, rounded corners=3pt, inner sep=4pt},
  % ---- state-sequence frames: show an operation happen step by step ----
  % (rotations, BST insert/delete, binary-lifting table fill, HLD splits,
  %  lazy push-down, centroid cuts). Draw each snapshot with the tree styles
  %  above, then wrap it in a frame and link frames with snap arrows.
  frame/.style={rectangle, rounded corners=2pt, draw=darkgray!55, fill=white,
    inner sep=6pt},
  framehl/.style={frame, draw=accentorange, fill=accentorange!5},
  framecap/.style={font=\footnotesize\itshape, text=darkgray, align=center},
  stepno/.style={circle, fill=primaryblue, text=white, inner sep=2pt,
    minimum size=6mm, font=\bfseries\scriptsize},
  snap/.style={->, very thick, accentorange},
}

\pgfplotsset{
  bookplot/.style={
    width=0.8\linewidth, height=5.5cm, grid=major, grid style={lightgray},
    axis lines=left, tick label style={font=\footnotesize},
    label style={font=\small}, legend style={font=\footnotesize, draw=none, fill=none},
    every axis plot/.append style={thick}},
}

% forest trees (binary trees, general trees, recursion trees)
\forestset{
  booktree/.style={for tree={circle, draw=primaryblue, fill=primaryblue!8, thick,
    minimum size=7mm, inner sep=1pt, font=\small, l sep=6mm, s sep=4mm, edge={thick, darkgray}}},
  bintree/.style={for tree={circle, draw=primaryblue, fill=primaryblue!8, thick,
    minimum size=8mm, inner sep=0pt, font=\small, l=11mm, s sep=7mm,
    edge={thick, darkgray}, calign=fixed edge angles,
    calign primary angle=-40, calign secondary angle=40}},
  genetree/.style={for tree={circle, draw=primaryblue, fill=primaryblue!8, thick,
    minimum size=7mm, inner sep=0pt, font=\small, l sep=8mm, s sep=5mm,
    edge={thick, darkgray}}},
  valbintree/.style={for tree={rectangle, rounded corners=1pt, draw=darkgray,
    fill=white, thick, minimum size=7mm, inner sep=1pt, font=\footnotesize,
    l=10mm, s sep=5mm, edge={thick, darkgray}}},
  % small forest tree for one frame of a state sequence (side-by-side snapshots)
  snaptree/.style={for tree={circle, draw=primaryblue, fill=primaryblue!8, thick,
    minimum size=6mm, inner sep=0pt, font=\footnotesize, l=8mm, s sep=4mm,
    edge={thick, darkgray}}},
}
